% ================================================================
% The Retrieve and Connect version of latex/starters/targeted/KS3/statistics/developingAveragesStarters.tex
%
% This is the same deck as the one above, made by renaming the titles in it.
% Every slide that was a numbered starter is now titled
% "Retrieve and Connect", and the word is renamed wherever else a title
% used it. Nothing outside a title has been changed.
%
%   made from           latex/starters/targeted/KS3/statistics/developingAveragesStarters.tex
%   retitling rules     version 1
%
% Please don't edit this file. It is written again from the one it was
% made from whenever that changes, so an edit here would be lost - edit
% that file instead and this one follows.
% ================================================================
\documentclass{beamer}
\usepackage{tikz}
\usetikzlibrary{calc,arrows.meta,patterns}
\usepackage{xcolor}
\usepackage{multicol}
\usepackage{amsmath}

% ================================================================
% Answer-overlay helpers
% ================================================================
\newcommand{\ablank}[1]{%
  \alt<2>{\textcolor{red}{#1}}{\underline{\phantom{#1}}}%
}
\newcommand{\ashow}[1]{\uncover<2->{\textcolor{red}{\small #1}}}
\newcommand{\ashowq}[1]{\alt<2->{\textcolor{red}{\small #1}}{?}}

% Answers drawn straight on to a TikZ diagram, revealed on overlay 2 like \ashow.
% Space is reserved on overlay 1, so the diagram does not move between slides.
\usetikzlibrary{overlay-beamer-styles}
\tikzset{
  ans/.style     = {red, thick, visible on=<2->},
  ansfill/.style = {red!15, visible on=<2->},
  anslab/.style  = {red, font=\tiny, visible on=<2->},
}

\title{Averages: Mean \& Median}
\author{}
\date{}

% Lego tower: \legotower{x-offset}{height}{colour}
% Each brick is 0.65 wide x 0.55 tall, with two studs
\newcommand{\legotower}[3]{%
  \foreach \i in {1,...,#2}{%
    \draw[fill=#3!55, draw=#3!80!black, line width=0.35pt]
      (#1, {(\i-1)*0.55}) rectangle (#1+0.65, {\i*0.55});
    \draw[fill=#3!80, draw=#3!80!black, line width=0.2pt]
      (#1+0.15,{(\i-1)*0.55+0.38}) circle (0.08);
    \draw[fill=#3!80, draw=#3!80!black, line width=0.2pt]
      (#1+0.44,{(\i-1)*0.55+0.38}) circle (0.08);
  }%
}

\begin{document}


\begin{frame}{Retrieve and Connect}
\vspace{-0.8em}

\footnotesize
\begin{enumerate}\itemsep2pt
\begin{multicols}{2}
  \item $6 + 9 = \ashowq{15}$
  \item $24 \div 4 = \ashowq{6}$
  \item $7 + 13 + 5 = \ashowq{25}$
  \item $30 \div 5 = \ashowq{6}$
  \end{multicols}
  \item How do we work out the mean? \ashow{Add the values and divide by how many there are.}
  \item How do we work out the median? \ashow{Order the data and find the middle value.}
  \item Find the mean of: $3,\ 7,\ 5,\ 9$ \ashow{$=6$}
  \item Find the median of: $4,\ 9,\ 2,\ 7,\ 8$ \ashow{$=7$}
  \item Find the mean of: $\frac{1}{2},\ \frac{3}{2},\ 2,\ 3$ \ashow{$=\frac{7}{4}$}
  \item Find the median of: $-4,\ 2,\ -1,\ -7,\ 1$ \ashow{$=-1$}
  \item The mean of $x,\ x+4,\ x+8,\ x+12$ is $15$, what is the total of all the values? \ashow{$60$}
  \item Find the value of $x$ in the previous question \ashow{$x=9$}
\end{enumerate}

\end{frame}


\begin{frame}{Retrieve and Connect -- Solutions}
\vspace{-0.8em}

\footnotesize
\begin{enumerate}\itemsep2pt
\begin{multicols}{2}
  \item $6 + 9 = \mathord{?}$ \textcolor{red}{$15$}
  \item $24 \div 4 = \mathord{?}$ \textcolor{red}{$6$}
  \item $7 + 13 + 5 = \mathord{?}$ \textcolor{red}{$25$}
  \item $30 \div 5 = \mathord{?}$ \textcolor{red}{$6$}
\end{multicols}
  \item How do we work out the mean?  
        \textcolor{red}{Add the values and divide by how many there are}
  \item How do we work out the median?  
        \textcolor{red}{Order the data and find the middle value}
  \item Find the mean of: $3,\ 7,\ 5,\ 9$  
        \textcolor{red}{$=6$}
  \item Find the median of: $4,\ 9,\ 2,\ 7,\ 8$  
        \textcolor{red}{$=7$}
  \item Find the mean of: $\frac{1}{2},\ \frac{3}{2},\ 2,\ 3$  
        \textcolor{red}{$=\frac{7}{4}$}
  \item Find the median of: $-4,\ 2,\ -1,\ -7,\ 1$  
        \textcolor{red}{$=-1$}
  \item The mean of $x,\ x+4,\ x+8,\ x+12$ is $15$, what is the total?  
        \textcolor{red}{$60$}
  \item Find the value of $x$  
        \textcolor{red}{$x=9$}
\end{enumerate}
\end{frame}
%------------------------------------------------
% STARTER 2 — Mean and Median
%------------------------------------------------
\begin{frame}{Retrieve and Connect}
\vspace{-0.8em}

\footnotesize
\begin{enumerate}\itemsep2pt
\begin{multicols}{2}
  \item $3 + 7 + 10 = \ashowq{20}$
  \item $48 \div 6 = \ashowq{8}$
  \item $\frac{1}{2} + \frac{1}{4} = \ashowq{\frac{3}{4}}$
  \item $28 \div 4 = \ashowq{7}$
  \end{multicols}
  \item Which average requires you to order the data first? \ashow{Median}
  \item Find the median of: $7,\ 2,\ 9,\ 4,\ 5$ \ashow{$=5$}
  \item Find the median of: $3,\ 8,\ 1,\ 6,\ 4,\ 10$ \ashow{$=5$}
  \item Find the mean and median of: $2,\ 5,\ 3,\ 8,\ 7$.  
        Which is larger? \ashow{Mean $=5$, median $=5$ (equal).}
  \item Find the mean and median of: $-5,\ 1,\ -2,\ 4,\ 7$ \ashow{Mean $=1$, median $=1$.}
  \item The mean of $a,\ a+3,\ a+6$ is $6$.
        \begin{itemize}\footnotesize
        \item What is the median in terms of $a$? \ashow{$a+3$}
          \item What is the total of the data (use the mean) \ashow{$18$}
          \item Find $a$ \ashow{$a=3$}
        \end{itemize}
\end{enumerate}

\end{frame}

\begin{frame}{Retrieve and Connect -- Solutions}
\vspace{-0.8em}

\footnotesize
\begin{enumerate}\itemsep2pt
\begin{multicols}{2}
  \item $3 + 7 + 10 = \mathord{?}$ \textcolor{red}{$20$}
  \item $48 \div 6 = \mathord{?}$ \textcolor{red}{$8$}
  \item $\frac{1}{2} + \frac{1}{4} = \mathord{?}$ \textcolor{red}{$\frac{3}{4}$}
  \item $28 \div 4 = \mathord{?}$ \textcolor{red}{$7$}
\end{multicols}
  \item Which average requires ordering first?  
        \textcolor{red}{Median}
  \item Median of $7,\ 2,\ 9,\ 4,\ 5$  
        \textcolor{red}{$=5$}
  \item Median of $3,\ 8,\ 1,\ 6,\ 4,\ 10$  
        \textcolor{red}{$=5$}
  \item Mean and median of $2,\ 5,\ 3,\ 8,\ 7$  
        \textcolor{red}{Mean $=5,\ \text{median }=5$}
  \item Mean and median of $-5,\ 1,\ -2,\ 4,\ 7$  
        \textcolor{red}{Mean $=1,\ \text{median }=1$}
  \item Mean of $a,\ a+3,\ a+6$ is $6$
        \begin{itemize}\footnotesize
          \item Median \textcolor{red}{$=a+3$}
          \item Total \textcolor{red}{$=18$}
          \item $a=$ \textcolor{red}{$3$}
        \end{itemize}
\end{enumerate}
\end{frame}

%------------------------------------------------
% STARTER 3 — Mean, Median, Mode
%------------------------------------------------
\begin{frame}{Retrieve and Connect}
\vspace{-0.8em}

\footnotesize
\begin{enumerate}\itemsep2pt
\begin{multicols}{2}
  \item $4 + 11 = \ashowq{15}$
  \item $40 \div 8 = \ashowq{5}$
  \item $\frac{3}{4} - \frac{1}{4} = \ashowq{\frac{1}{2}}$
  \item $72 \div 9 = \ashowq{8}$
  \end{multicols}
  \item A shoe shop owner says,  
        ``The most popular size last week was size 6.''  
    Which average is she using? \ashow{Mode}
  \item Find the mode of: $3,\ 7,\ 3,\ 9,\ 5,\ 3,\ 7$ \ashow{$3$}
  \item Find the mode of: $2,\ 5,\ 8,\ 5,\ 3,\ 8,\ 1$ \ashow{$5$ and $8$}
  \item Find the mean and median of: $4,\ 7,\ 2,\ 9,\ 3$ \ashow{Mean $=5$, median $=4$.}
  \item Find the mean, median and mode of: $5,\ 3,\ 7,\ 3,\ 5,\ 5,\ 9$ \ashow{Mean $=\frac{37}{7}$, median $=5$, mode $=5$.}
  \item The data set is $\{4,\ 7,\ 4,\ x,\ 7,\ 3\}$  
        The \textbf{only} mode is $4$.  
        Find $x$ and explain your reasoning. \ashow{$x=4$; then $4$ occurs three times.}
\end{enumerate}

\end{frame}


\begin{frame}{Retrieve and Connect -- Solutions}
\vspace{-0.8em}

\footnotesize
\begin{enumerate}\itemsep2pt
\begin{multicols}{2}
  \item $4 + 11 = \mathord{?}$ \textcolor{red}{$15$}
  \item $40 \div 8 = \mathord{?}$ \textcolor{red}{$5$}
  \item $\frac{3}{4} - \frac{1}{4} = \mathord{?}$ \textcolor{red}{$\frac{1}{2}$}
  \item $72 \div 9 = \mathord{?}$ \textcolor{red}{$8$}
\end{multicols}
  \item Which average is used?  
        \textcolor{red}{Mode}
  \item Mode of $3,\ 7,\ 3,\ 9,\ 5,\ 3,\ 7$  
        \textcolor{red}{$3$}
  \item Mode of $2,\ 5,\ 8,\ 5,\ 3,\ 8,\ 1$  
        \textcolor{red}{$5\ \text{and}\ 8$}
  \item Mean and median of $4,\ 7,\ 2,\ 9,\ 3$  
        \textcolor{red}{Mean $=5,\ \text{median }=4$}
  \item Mean, median and mode of $5,\ 3,\ 7,\ 3,\ 5,\ 5,\ 9$  
        \textcolor{red}{Mean $=\frac{37}{7}$, median $=5$, mode $=5$}
  \item Only mode is $4$  
        \textcolor{red}{$x=4$}
\end{enumerate}
\end{frame}
%------------------------------------------------
% STARTER 4 — Consolidation
%------------------------------------------------
\begin{frame}{Retrieve and Connect}
\vspace{-0.8em}

\footnotesize
\begin{enumerate}\itemsep2pt
\begin{multicols}{2}
  \item $3+4 = \ashowq{7}$
  \item $4 \times 8 = \ashowq{32}$
  \item $-3 + 8 + (-5) = \ashowq{0}$
  \item $\frac{2}{3} + \frac{1}{6} = \ashowq{\frac{5}{6}}$
  \item $-12 \div 4 = \ashowq{-3}$
  \item $\frac{3}{4} \times 8 = \ashowq{6}$
  \end{multicols}
  \item A class votes for their favourite colour.  
        Can you find the mean?  
        Which average should you use instead, and why?  
        \ashow{No; use the mode, because the data are categories.}
  \item Find the mean, median and mode of:  
        $-3,\ 1,\ -3,\ 4,\ -1,\ 2,\ -3$ \ashow{Mean $=-\frac{3}{7}$, median $=-1$, mode $=-3$.}
  \item Find the median of:  
        $\frac{1}{4},\ \frac{3}{4},\ \frac{1}{2},\ \frac{1}{8},\ \frac{5}{8}$ \ashow{$\frac{1}{2}$}
  \item The ages in a youth club are:  
        $11,\ 12,\ 12,\ 13,\ 13,\ 13,\ 14,\ 42$
        Which average would you not use for this example, and why? \ashow{Mean, because $42$ is an outlier.}
  \item The mean of $\{x,\ 3,\ x,\ 7,\ 5\}$ is $6$.  
        Find $x$ and state the mode. \ashow{$x=6$, mode $=6$}
  \item Five numbers have mean $8$, median $7$ and unique mode of $5$.  
        Two numbers are $5$ and $5$.  
        Find a possible set. \ashow{$\{5,5,6,7,17\}$}
\end{enumerate}

\end{frame}

\begin{frame}{Retrieve and Connect -- Solutions}
\vspace{-0.8em}

\footnotesize
\begin{enumerate}\itemsep2pt
\begin{multicols}{2}
  \item $3+4 = \mathord{?}$ \textcolor{red}{$7$}
  \item $4 \times 8 = \mathord{?}$ \textcolor{red}{$32$}
  \item $-3 + 8 + (-5) = \mathord{?}$ \textcolor{red}{$0$}
  \item $\frac{2}{3} + \frac{1}{6} = \mathord{?}$ \textcolor{red}{$\frac{5}{6}$}
  \item $-12 \div 4 = \mathord{?}$ \textcolor{red}{$-3$}
  \item $\frac{3}{4} \times 8 = \mathord{?}$ \textcolor{red}{$6$}
\end{multicols}
  \item Can you find the mean?  
        \textcolor{red}{No — data is categorical}
  \item Mean, median and mode  
        \textcolor{red}{Mean $=-\frac{3}{7}$, median $=-1$, mode $=-3$}
  \item Median  
        \textcolor{red}{$=\frac{1}{2}$}
  \item Which average not to use and why?  
        \textcolor{red}{Mean, because 42 is an outlier}
  \item $x=?$ and mode  
        \textcolor{red}{$x=6$, mode $=6$}
  \item Possible set  
        \textcolor{red}{$\{5,5,6,7,17\}$}
\end{enumerate}
\end{frame}

%------------------------------------------------
% STARTER 5 — Challenge
%------------------------------------------------
\begin{frame}{Retrieve and Connect}
\vspace{-0.8em}

\footnotesize
\begin{enumerate}\itemsep2pt
\begin{multicols}{2}
\item $3+8 = \ashowq{11}$
\item $120 \div 10 = \ashowq{12}$
  \item $-6 + (-2) + 5 = \ashowq{-3}$
  \item Order from smallest to largest:  
        $\frac{1}{3},\ \frac{5}{6},\ \frac{1}{2},\ \frac{2}{3}$ \ashow{$\frac{1}{3},\ \frac{1}{2},\ \frac{2}{3},\ \frac{5}{6}$}
  \item $3\frac{1}{2} + 2\frac{3}{4} = \ashowq{6\frac{1}{4}}$
  \item $-20 \div 4 = \ashowq{-5}$
  \end{multicols}

  \item Find the mean, median and mode of:  
        $-6,\ -2,\ 4,\ -3,\ 7,\ 0,\ -2$ \ashow{Mean $=-\frac{2}{7}$, median $=-2$, mode $=-2$.}

  \item True or False?  
        ``Adding 10 to every value increases the mean by 10 but leaves the median unchanged.''  
        \ashow{True}

  \item The seven values  
        $n-9,\ n-6,\ n-3,\ n,\ n+3,\ n+6,\ n+9$  
        have a mean of $4$.  
        State the median \textbf{without calculation}. \ashow{Median $=4$}
  \item A set of five positive integers has  
        mean = median = mode = $6$.
        \begin{itemize}\footnotesize
          \item What is the smallest possible range? \ashow{$1$}
          \item What is the largest possible range? \ashow{$12$}
        \end{itemize}
\end{enumerate}

\end{frame}

\begin{frame}{Retrieve and Connect -- Solutions}
\vspace{-0.8em}

\footnotesize
\begin{enumerate}\itemsep2pt
\begin{multicols}{2}
  \item $3+8 = \mathord{?}$ \textcolor{red}{$11$}
  \item $120 \div 10$ \textcolor{red}{$12$}
  \item $-6 + (-2) + 5$ \textcolor{red}{$-3$}
  \item Order  
        \textcolor{red}{$\frac13,\ \frac12,\ \frac23,\ \frac56$}
  \item $3\frac12 + 2\frac34$ \textcolor{red}{$6\frac14$}
  \item $-20 \div 4$ \textcolor{red}{$-5$}
\end{multicols}
  \item Mean, median, mode  
        \textcolor{red}{Mean $=-\frac{2}{7}$, median $=-2$, mode $=-2$}
  \item True or False?  
        \textcolor{red}{True}
  \item Median  
        \textcolor{red}{$4$}
  \item Range  
        \textcolor{red}{Smallest $=1$, largest $=12$}
\end{enumerate}
\end{frame}

\end{document}